\section{An exponential encoding lower bound}\label{sec:lower}
All encoding statements use explicit rational coefficients and ordinary
binary representations of their numerators and positive denominators.
A sparse encoding records variable indices as well as nonzero coefficients.
The norm and the relaxed equality are fixed as displayed.

For $n\ge1$, let
\begin{equation}\label{eq:chain}
 A_n=\{a\in[0,1]^n:a_1\ge1/4,\ a_{i+1}\ge a_i^2\ (1\le i<n)\},
 \qquad \delta_n=2^{-2^n}.
\end{equation}
Induction shows $a_i\ge2^{-2^i}$, with all bounds attained by
$\bar a_i=2^{-2^i}$. Define
\begin{equation}\label{eq:onebinary}
 X_n=\{(a,y,q):a\in A_n,\ -1\le y\le1,\ q\in\{0,1\},\
                  y\ge a_n-2(1-q)\}.
\end{equation}
We minimize $-q$ subject to $(a,y,q)\in X_n$ and $y=0$.
In \eqref{eq:model}, take $f=-q$, $r=y$, and $\norm r=|y|$.

\begin{theorem}\label{thm:lower}
The problem defined using $X_n$ in \eqref{eq:onebinary} has value $v_n=0$, and for every $\rho\ge0$,
\begin{equation}\label{eq:dualformula}
 D_{n,\rho}=\min\left\{0,\frac{2\rho\delta_n-1}{1+\delta_n}\right\}.
\end{equation}
Its dual supremum is attained and its least dual-value exact penalty is
\begin{equation}\label{eq:lowerthreshold}
 \rho_{*,n}=\frac1{2\delta_n}=2^{2^n-1}.
\end{equation}
The instance has one binary variable, $n+1$ continuous variables, a linear
objective, and convex quadratic native constraints with coefficients from
a fixed finite set. Its sparse encoding length is $O(n\log(n+1))$, whereas
every rational exact penalty needs at least $2^n$ numerator bits.
Every equality-feasible integer slice and every native integer slice have
strict points whose continuous inequality slacks are at least $1/8$.
\end{theorem}
\begin{proof}
Projection onto $(q,y)$ gives exactly
\begin{equation}\label{eq:projection}
 (\{0\}\times[-1,1])\ \cup\ (\{1\}\times[\delta_n,1]).
\end{equation}
The inclusion follows from the chain; the reverse follows by setting
$a=\bar a$. At $q=0$ the lower bound $\delta_n-2$ is below $-1$.
Consequently $y=0$ forces $q=0$ and $v_n=0$.

For any multiplier, the points $(q,y)=(0,-1)$ and $(1,\delta_n)$
give values $\rho-\lambda$ and $-1+(\rho+\lambda)\delta_n$.
Their weighted average, with weights $\delta_n/(1+\delta_n)$ and
$1/(1+\delta_n)$, is $(2\rho\delta_n-1)/(1+\delta_n)$.
The point $(0,0)$ gives zero. These witnesses prove the upper bound
in \eqref{eq:dualformula} for every multiplier.

For $0\le\rho\le1/(2\delta_n)$, take
\[
 \lambda_\rho=\frac{1+\rho(1-\delta_n)}{1+\delta_n}\ge\rho.
\]
On the first interval of \eqref{eq:projection}, the minimum occurs at
$y=-1$; on the second it occurs at $y=\delta_n$. Their values coincide
with the proposed dual value. For $\rho\ge1/(2\delta_n)$, take
$\lambda=\rho$. The penalized objective is nonnegative on both intervals
and zero at $(0,0)$. This proves the formula and attainment.

Each inequality $a_i^2-a_{i+1}\le0$ is convex; all other constraints
are affine. Across all constraints there are $O(n)$ nonzero coefficients,
each with variable indices of $O(\log(n+1))$ bits.
The integer $2^{2^n-1}$ has exactly $2^n$ bits.
If a rational $p/d$ with positive integers $p,d$ is at least this value,
then $p\ge2^{2^n-1}$ as well.

For the equality-feasible slice $q=0$, set $a_i=3/8$ and $y=0$.
For the native slice $q=1$, set $a_i=3/8$ and $y=1/2$.
The seed inequality has slack $1/8$, each chain inequality has slack
$15/64$, and the variable bounds and linking inequality have slack at
least $1/8$. The native $q=0$ slice uses the first point. These are
uniform strict points, although the equality-infeasible $q=1$ slice
comes within $\delta_n$ of $y=0$.
\end{proof}
The construction satisfies Assumptions~1--5 of
\citet{LefebvreSchmidt2025}, so their Theorem~14 ensures a finite exact
penalty at every fixed multiplier. Theorem~\ref{thm:lower} concerns the
encoding length of such penalties, even after optimizing the multiplier.

At $\rho=\rho_{*,n}$, Proposition~\ref{prop:scalar} gives the unique exact
multiplier $\lambda=\rho_{*,n}$: here $A_+=1/\delta_n$ and $A_-=0$.
Both witnesses in the proof then tie with feasible optima.
For every $\rho>\rho_{*,n}$, the fixed multiplier
$\lambda=\rho_{*,n}$ makes all infeasible objective values positive,
so the minimizer set is exact. The zero-multiplier value threshold is
$1/\delta_n$, twice the optimized threshold.

\begin{corollary}[Fixed additive accuracy]\label{cor:accuracy}
For $0\le\epsilon<1$, the condition $v_n-D_{n,\rho}\le\epsilon$ is
equivalent to
\begin{equation}\label{eq:accuracy}
 \rho\ge\max\left\{0,\frac{1-\epsilon(1+\delta_n)}{2\delta_n}\right\}.
\end{equation}
For each fixed $\epsilon<1$ and all sufficiently large $n$, such a rational
penalty needs at least $2^n-O_\epsilon(1)$ numerator bits.
\end{corollary}
\begin{proof}
The first claim follows from \eqref{eq:dualformula}.
For large $n$, $\epsilon\delta_n\le(1-\epsilon)/2$, so the right-hand side
is at least $(1-\epsilon)/(4\delta_n)$. Taking binary logarithms gives
the bit bound.
\end{proof}
In particular, accuracy $\epsilon=1/2$ requires
$\rho\ge(1-\delta_n)/(4\delta_n)$. The objective range is one, so this
is a fixed accuracy relative to that normalization. The lower bound is
exponential in continuous dimension and superpolynomial in input length;
it is not a claim of an exponential bound in the full sparse input length.
Nor is it a running-time lower bound for solving the primal, whose value
is immediate from \eqref{eq:projection}.

\subsection{A symmetric variant and the role of the penalty function}
A two-binary companion makes multiplier cancellation symmetric.
Replace the final inequality in \eqref{eq:onebinary} by
\begin{equation}\label{eq:symmetric}
 y\ge a_n-(1-q)-2(1-b),\qquad
 y\le-a_n+(1-q)+2b,\qquad b\in\{0,1\}.
\end{equation}
Keep the same objective and relaxed equality. Projection now gives
\[
 (\{0\}\times[-1,1])\ \cup\
 (\{1\}\times([-1,-\delta_n]\cup[\delta_n,1])).
\]
Indeed, at $a=\bar a$ the two $q=0$ intervals are
$[-1,1-\delta_n]$ and $[\delta_n-1,1]$, which cover $[-1,1]$;
at $q=1$ they are the two outer intervals. If $c=\rho-|\lambda|$,
minimizing on these intervals yields
\begin{equation}\label{eq:symdual}
 L_{n,\rho}^{\rm sym}(\lambda)=\min\{0,-1+c\delta_n,-1+c\},\qquad
 D_{n,\rho}^{\rm sym}=\min\{0,-1+\rho\delta_n\}.
\end{equation}
For the first formula, the $q=0$ branch has value $\min(0,c)$ and
is dominated by one of the three displayed terms. The second follows
because $c\le\rho$ and $\lambda=0$ attains the bound.
Thus the least exact penalty is $\delta_n^{-1}$, and gap at most
$\epsilon<1$ requires exactly $\rho\ge(1-\epsilon)\delta_n^{-1}$.
Any strictly larger penalty with $\lambda=0$ is minimizer-set exact.
Strict points again use $a_i=3/8$, with $y=0$ for $q=0$ and
$y=(2b-1)/2$ for $q=1$; their minimum continuous inequality slack is $1/8$.

The fixed norm is essential. Let $p_n=2^{-n}$, a rational exponent with
$O(n)$ bits. Then $\delta_n^{p_n}=1/2$. For either construction,
$f+2|y|^{p_n}$ has minimum zero: at $q=0$ it is nonnegative and at
$q=1$ it is at least $-1+2\delta_n^{p_n}=0$.
Any coefficient larger than two gives minimizer-set exactness.
This instance-dependent augmentation is not Lipschitz at zero and is
not a norm. It changes the encoding question and the continuous
minimization problem. Likewise, succinct exponent expressions or a
rescaled linking equality change the representation or formulation to
which Theorem~\ref{thm:lower} applies.
